\section{Geodesics, Busemann functions, and structural input}
\label{sec:preliminaries}

\subsection{The metric and its rays}
Let \(\mathcal E\) be the unoriented nearest-neighbor edges of \(\Z^2\).
We work on
\[
 \Omega=[0,\infty)^{\mathcal E},\qquad \PP=G^{\mathcal E},
\]
initially with the product Borel sigma-algebra and subsequently with its
completion. Throughout the paper \(G\) is nonatomic and satisfies the
moment assumption of Theorem~\ref{thm:main}.
For a finite nearest-neighbor walk \(\gamma\), its cost \(T(\gamma)\)
is the sum of its edge weights with traversal multiplicity. Put
\[
 T(x,y)=\inf_{\gamma:x\longrightarrow y}T(\gamma).
\]
A finite \emph{geodesic} realizes this infimum. A \emph{ray}
\(p=(p_0,p_1,\ldots)\) is a simple infinite nearest-neighbor path every
finite segment of which is a geodesic. A \emph{bigeodesic} has the same
property and is indexed by \(\Z\).
Two rays \emph{coalesce} if, after deleting finite initial portions,
their vertex sequences agree.
We draw lattice edges as straight segments; intersections of such paths
are lattice contacts.

\begin{lemma}[Finite minimizers]\label{lem:finite-minimizers}
Almost surely every pair of vertices has a unique finite geodesic.
\end{lemma}
\begin{proof}
Every edge weight is strictly positive almost surely. By dominated
convergence, choose \(c>0\) such that
\(a=4\EE e^{-ct_e}<1\).
For a fixed vertex \(x\) and a fixed \(K<\infty\), the probability that
some simple path of \(j\) edges from \(x\) has cost at most \(K\) is
at most \(e^{cK}a^j\). The sum over \(j\) is finite.
Consequently only finitely many simple paths from \(x\) have cost at
most \(K\), almost surely. Take the countable intersection over vertices
and positive integer \(K\).
Every finite walk can be shortened to a simple path by erasing loops.
Comparing with any fixed path between two vertices now shows that the
infimum is attained among a finite collection of paths.

Two different simple paths with the same endpoints have an edge in their
symmetric difference. Conditional on all other weights, equality of their
costs prescribes one value of the independent nonatomic weight on that
edge and thus has probability zero. There are only countably many pairs
of finite paths.
\end{proof}

The raywise monotonicity argument below goes back to
Hoffman~\cite{Hoffman2005}; see also
\cite[Lemma~4.1 and the discussion after Proposition~4.5]{AhlbergHoffman}.

\begin{lemma}[Elementary Busemann properties]\label{lem:busemann-basic}
For every ray \(p\) and all vertices \(x,y\), the limit
\begin{equation}\label{eq:busemann-definition}
 B_p(x,y)=\lim_{j\to\infty}\bigl(T(x,p_j)-T(y,p_j)\bigr)
\end{equation}
exists. It is additive and antisymmetric, satisfies
\(\lvert B_p(x,y)\rvert\le T(x,y)\), and is calibrated on the ray:
\[
 B_p(p_i,p_j)=T(p_i,p_j)\qquad(i\le j).
\]
Deleting a finite initial portion of a ray does not change its Busemann
function. Coalescing rays have the same Busemann function. For each
fixed vertex \(x\),
\begin{equation}\label{eq:calibration-defect}
 0\le T(x,p_j)-B_p(x,p_j)\longrightarrow0.
\end{equation}
\end{lemma}
\begin{proof}
The sequence
\[
 a_j(x)=T(x,p_j)-T(p_0,p_j)
\]
is decreasing, by the triangle inequality and calibration of the
segments of \(p\), and is bounded below by \(-T(p_0,x)\).
Write \(a_\infty(x)\) for its limit. Then
\(B_p(x,y)=a_\infty(x)-a_\infty(y)\).
Additivity, antisymmetry, and the metric bound follow immediately.
For \(j\ge i\) and \(k\ge j\), the geodesic property gives
\(T(p_i,p_k)-T(p_j,p_k)=T(p_i,p_j)\), proving calibration.
Tail invariance and the coalescence assertion follow from the defining
limit. Finally,
\[
 B_p(x,p_j)=B_p(x,p_0)+T(p_0,p_j)
           =a_\infty(x)+T(p_0,p_j),
\]
so the left side of \eqref{eq:calibration-defect} is
\(a_j(x)-a_\infty(x)\).
\end{proof}

\subsection{Structural input}
For a norm \(m\) on \(\R^2\), let
\[
 \mathcal S=\partial\{\rho\in(\R^2)^*:\rho(v)\le m(v)
                                      \text{ for every }v\in\R^2\},
 \qquad
 A(\rho)=\{d\in S^1:m(d)=\rho(d)\}.
\]
Here \(|\cdot|\) and \(S^1\) refer to the Euclidean norm. The set
\(\mathcal S\) is the dual unit circle, and \(A(\rho)\) is the set of
contact directions of \(\rho\).
The limiting directions of a ray are the accumulation points of
\((p_j-p_0)/|p_j-p_0|\).
Since a simple lattice ray eventually leaves each finite set, this is
a nonempty compact subset of \(S^1\).

The structural input below combines the shape theorem with the
Ahlberg--Hoffman theory. The boundary comparison in
Appendix~\ref{sec:boundary-comparison} is part of this input: it replaces
a boundary-expectation step without adding a moment assumption.

\begin{theorem}[Ahlberg--Hoffman and the shape theorem]
\label{thm:structure}
There are a deterministic norm \(m\) on \(\R^2\) and a deterministic
closed subset \(C\subseteq\mathcal S\) with the following properties.
\begin{enumerate}[label=\textup{(\roman*)}]
\item Almost surely,
\[
 T(0,x)=m(x)+o(|x|)\qquad (|x|\to\infty).
\]
\item Almost surely, simultaneously for all rays \(p\) from all lattice
vertices, there is a unique \(\rho_p\in C\) such that
\begin{equation}\label{eq:busemann-linearity}
 B_p(0,x)=\rho_p(x)+o(|x|)\qquad (|x|\to\infty).
\end{equation}
The limiting directions of \(p\) lie in \(A(\rho_p)\).
At each vertex the set of realized functionals is \(C\).
\item For each fixed \(\rho\in C\), almost surely there is exactly one
ray \(g_\rho(x)\) with functional \(\rho\) from each \(x\in\Z^2\).
These rays coalesce. Their family is \emph{backwards finite}: for each
vertex \(z\), the set of \(x\) for which \(g_\rho(x)\) contains \(z\)
is finite.
\end{enumerate}
The norm \(m\) and the set \(C\) have the rotation and reflection
symmetries of the square lattice.
\end{theorem}

The exact source statements are
\cite[Assumption A1, Theorems 2.5--2.6 and 10.8--10.11,
and Proposition 4.5]{AhlbergHoffman}.
Its Assumption A1 is the iid continuous-marginal setting with the
second moment of the minimum of four incident edge weights.
A nonatomic Borel law has a continuous distribution function, so
singular continuous laws are included.
The iid shape theorem is due to Cox and Durrett~\cite{CoxDurrett1981}
and is recalled there in Section~3.1.
Theorem~10.8 supplies a random coalescing geodesic for each fixed
member of \(C\); Theorem~10.11 and its proof identify its translates
as the unique rays of that functional, and Proposition~4.5 gives
backward finiteness. Countability of the lattice permits all starting
vertices in the displayed assertions.
The separate boundary-expectation assertion of their Lemma~4.6 is not
needed: the appendix replaces its use in Proposition~4.4 by a finite
stationary correction and bounded truncation. The geometric
existence and coalescence part of that lemma is unchanged.
The symmetry of \(C\) also follows from its deterministic realization
as the set of ray functionals.

The distinction between parts \textup{(ii)} and \textup{(iii)} is
essential. Part \textup{(ii)} holds for every ray in a single typical
environment. In part \textup{(iii)}, the null set may depend on the
\emph{fixed} functional. We will not intersect those events over an
uncountable collection of functionals.

For a fixed \(\rho\in C\), the coalescing family has a common Busemann
function, denoted \(B_\rho\). It is defined on the probability-one
event for that fixed functional. For any countable deterministic
collection of such functionals we use their common probability-one
event.
We call \(\rho_p\) the \emph{label} of \(p\).

\begin{corollary}[Exclusion of deterministic labels]
\label{cor:fixed-label}
For each fixed \(\rho\in C\), almost surely no bigeodesic has an end
with label \(\rho\). The same holds simultaneously for every member
of any countable deterministic subset of \(C\).
\end{corollary}
\begin{proof}
Suppose a bigeodesic has a forward end of label \(\rho\).
The forward ray starting at any of its backward vertices has that
same label by tail invariance. Fixed-label uniqueness identifies all
of these rays with \(g_\rho(x)\).
The middle vertex would belong to this family from infinitely many
starting vertices, contrary to backward finiteness.
The final assertion is a countable intersection.
\end{proof}

\subsection{Uniform estimates and finite changes}
\begin{lemma}[Uniform metric approximation]\label{lem:uniform-metric}
Almost surely,
\begin{equation}\label{eq:uniform-metric}
 T(x,y)=m(y-x)+o(|x|+|y|)
 \qquad (|x|+|y|\to\infty),
\end{equation}
uniformly over the two endpoints.
\end{lemma}
\begin{proof}
The extended shape theorem
\cite[Proposition 3.4]{AhlbergHoffman} states that for each
\(\eps>0\), almost surely there is an \(N<\infty\) such that
\[
 |T(z,z+v)-m(v)|\le\eps\max\{|z|,|v|\}
 \qquad (|z|\ge N,\ v\in\Z^2).
\]
Given \(x,y\), take \(z\) to be the endpoint with larger norm and \(v\)
the displacement to the other endpoint.
Then \(|z|\ge (|x|+|y|)/2\) and
\(\max\{|z|,|v|\}\le |x|+|y|\).
Symmetry of \(T\) and \(m\) proves the claim.
\end{proof}

\begin{lemma}[Finite modifications]\label{lem:finite-change}
Suppose two environments differ on a finite set \(F\) of edges, and put
\(K=\sum_{e\in F}|t'_e-t_e|\). Then
\[
 |T'(x,y)-T(x,y)|\le K\qquad(x,y\in\Z^2).
\]
If a path is a ray in both environments and has asymptotically linear
Busemann functions in both, its label is the same in the two environments.
The conclusion also applies when a common tail is a ray in both.
If \(t'_e\ge t_e\) for all \(e\) and a geodesic path uses only edges
whose weights are unchanged, it remains geodesic in the new metric.
\end{lemma}
\begin{proof}
Loop erasure allows the passage-time infimum to be taken over simple
paths. Each such path uses every changed edge at most once, so its
cost changes by at most \(K\). Taking infima gives the distance bound.
For a retained ray \(p\), passing to its defining Busemann limits gives
\[
 |B'_p(x,y)-B_p(x,y)|\le2K.
\]
Linear asymptotics along \(x=nv\), for two independent lattice vectors
\(v\), imply equality of the linear parts. Tail invariance handles
the stated extension. Under an increase, every competitor is at least
as expensive as before, whereas each segment of the retained path
has unchanged cost.
\end{proof}

Whenever a random environment has law absolutely continuous with respect
to \(\PP\), it inherits every specified \(\PP\)-almost-sure assertion.
We take common full-measure events for the countably many environments
and deterministic parameters used in each construction. In particular,
applying a simultaneous pathwise assertion in a finitely modified
environment does not require fixing the label of a subsequently selected
path.

\subsection{Measurability conventions}
\begin{lemma}[Path witnesses and selection]\label{lem:measurable}
The conditions that a specified path is a ray or a bigeodesic,
that specified rays coalesce, and that a ray's Busemann function is
asymptotically linear to a specified functional are Borel conditions
on the corresponding environment, path, and functional spaces.
Existence events obtained by projecting these conditions, with
additional Borel tests, are analytic and hence universally measurable.
On such existence events a witnessing path can be selected
completion-measurably.
The number of distinct labels with such witnesses is measurable in
the completed probability space, with the value \(+\infty\) allowed.
\end{lemma}
\begin{proof}
The spaces \(\Omega\), \((\Z^2)^{\mathbb N}\), and
\((\Z^2)^{\Z}\) are standard Borel spaces.
Nearest-neighbor adjacency and simplicity are countable coordinate
tests. Each finite segment minimizes exactly when its cost is no
greater than that of every finite competitor with the same endpoints;
there are countably many competitors. Equivalently, the passage
distance is a countable infimum of Borel finite-path costs.
Limits defining Busemann values are countable tests, as are
asymptotic linearity with rational tolerances and integer radii.
Coalescence is expressed by
\[
 \exists i,j\ge0\ \ \forall k\ge0:\ p_{i+k}=q_{j+k}.
\]
The standard analytic projection and Jankov--von Neumann selection
theorems give the existence and selection assertions; see
\cite[Theorems 29.7 and 29.9]{Kechris}.
To test that the label count is at least \(j\), project the Borel
condition for \(j\) path witnesses with pairwise distinct labels.
\end{proof}

We will also impose conditions on \emph{all} rays with labels in a
fixed chart from specified vertices. The failure of such a regularity
condition has one violating ray as an analytic witness.
The regularity event is therefore coanalytic and universally measurable.
It depends only on the environment and the specified vertices.
When combined with an existential path condition, we first project the
Borel witness relation and then intersect with this separate regularity
event. We do not project a relation containing an additional universal
path quantifier.

For a fixed law on a standard Borel space, a map measurable for its
completion and taking values in a standard Borel space has a Borel
version. Thus conditioning on a resulting environment fixes any path
chosen from that environment, up to a null set. Translated constructions
use the same selector on the relatively shifted input; their one-site
marginal laws are consequently identical. No stationarity or
independence of the resulting marks is needed beyond the explicitly
stated marginal conditions.
